\Page{6}{\textenglish{01 / DESIGN}}{לראות נפח לפני שיש דוגמה}{45}{53}

\FeatureIntro{PaperSteel}{נניח שמעצבת רואה על הלוח רצף של לוקים ארוכים וצרים, ורוצה לבדוק אם צללית רחבה תשבור אותו יפה. בתוכנה כמו \textenglish{CLO} היא יכולה לעבוד על קווי הגזרה, להציג את הבגד על דמות ממוחשבת ולראות איך גרסאות שונות יושבות זו לצד זו. אם הוגדרו תכונות בד מתאימות, ההדמיה יכולה לעזור לזהות איפה החומר נמתח ואילו אזורים כדאי לבדוק. על הדוגמה הפיזית אפשר לבחון את תחושת החומר ואת התנועה על הלובשת. כך כל שלב עונה על שאלה משלו: האם הרעיון עובד בתוך הקולקציה, ואיך הוא עובד כשמישהי באמת לובשת אותו. \SourceRef{2}}

\vspace{6mm}

{\fontsize{15}{23.2}\selectfont

החלון הדו־ממדי בתוכנת גזרות ממוחשבת, או \textenglish{CAD}, עוסק בחלקים שאפשר למדוד ולעדכן. כשמשנים שם קו מתאר, משתנה הגאומטריה של החלק. \textbf{דירוג מידות} מחיל על הגזרה כללים מתוכננים למידות נוספות. אפשר לשנות רוחב ואורך בשיעורים שונים, במקום להגדיל את הצורה כולה. בסימולציית \textbf{התאמה בתלת־ממד}, חלקי הגזרה מתחברים על דמות ממוחשבת. מפות ההתאמה של \textenglish{CLO} יכולות להצביע על אזורים שכדאי לבדוק. המידע על הגוף, החומר והגזרה קובע עד כמה אפשר לסמוך על המפות. \SourceRef{1}\SourceRef{2}

כלי מחולל נכנס בנקודה אחרת. הוא יכול להציע כיוונים חזותיים, לנסח טיוטת מפרט או, בכלי שמיועד לכך, להציע נקודת התחלה לגזרה. אפשר לתת לצוות לראות מוקדם יותר חלופה שבעבר הייתה נשארת כסקיצה. בשלב המעשי עדיין יש הבדל בין תמונה משכנעת, גאומטריה מסודרת ומוצר שיושב על הגוף כפי שהתכוונו. כשהסטודיו מתעד איזו מן השאלות הוא בודק בכל גרסה, הכלי הדיגיטלי יכול לכוון את העבודה על הדוגמה הבאה. \SourceRef{2}\SourceRef{3}

\par}
